\section{The Lemma Database}\label{sec:engine:db}

The database is a set of Dafny files. The entry point, \texttt{resources/cuvee/teq.dfy}, is injected with every
system: it declares the naturals, lists, traces and their combinators, counters, constant traces, and stores into and
reads of arrays, with the lemmas that every loop needs. Optional modules add concepts on demand
(\texttt{-include teq:<module>}): \emph{aggregation} (sums, counts, minima, maxima, parities), \emph{arithmetic}
(division, remainders, products, powers), \emph{arrays} (stores at distinct keys, conditional stores), \emph{counters}
(phases, flags, monotone and bounded values), \emph{search} (no occurrence in a prefix, the bisection of an interval),
\emph{sorted} (sortedness, permutations, partition points), \emph{strides} (strided stores), and \emph{unsigned}
(counters modulo $2^w$). Altogether there are 153 lemmas over 37 functions.

A lemma is used in one of three ways. A \emph{rewrite rule} (\texttt{\{:rw\}}) replaces its left-hand side, e.g.\ the
trace that adds $c$ in each step by the counter, and the last value of a counter by $i + nc$. A \emph{fold}
(\texttt{\{:fold\}}) is the definition of a function, whose body is rewritten to the function where it occurs; it names
a shape, such as $\stores$, $\mathsf{sums}$ or $\mathsf{bisection}$, and it is the preferred form over a declaration
with a separate equation, since the definition is then part of the semantics. A \emph{forward lemma}
(\texttt{\{:fw\}}) adds its instance wherever its trigger occurs among the assumptions
(Section~\ref{sec:engine:forward}), e.g.\ that a count is nonnegative or that a bisection keeps its lower end below the
cut of its guard. Rules are tried in the order of the files, so the order decides which abstraction a fact gets: rules
that name a trace come before rules that compute with named traces, and rules whose patterns subsume others come after
the specific ones.

The lemmas of the entry point and of the modules are proved in Lean~4 against a semantics of the combinators, traces as
nonempty lists and the combinators as folds, with guarded cases as conditionals that end in $\mathsf{unspecified}$; the
statements are generated from the Dafny files, so the proofs are about exactly the lemmas that the engine uses. Where a
lemma has arbitrary guards, it ends in a case for their complement (Section~\ref{sec:engine:cases}), so that complete
guards are not a premise. What the proofs still assume is stated as premises of the generated statements: reads of the
lists of a loop at its iterations, prefixes not longer than the list, lists of equal length, all of which the
construction of the systems guarantees. The proofs found two kinds of gaps, these implicit premises and an inconsistent
definition of the bisection (Section~\ref{sec:bisect}). The lemmas written for a single program, at the beginning of its
file, are not checked; Section~\ref{sec:principles:sound} discusses how to keep them sound.
